\documentclass[11pt]{article}
\usepackage{mathblatt}
% gesetzt von werkzeuge/setzer.py; Lösungen nach bau/bauregeln.md „Lösungen“.
\newcommand{\kennung}{3Y5}
\begin{document}
\pfheftstil
\fancyfoot[C]{{\footnotesize\color{mbgrau}\kennung\hfill\thepage}}
\pfheftkopf{Prozente als Anteile}{Klasse 7 \textperiodcentered{} Lösungen}

\begin{pfloesung}{}
\lz{1b)}{$8\,\%$}{$8$ von $100$ Kästchen $= \frac{8}{100}$}{}
\lz{c)}{$60\,\%$}{$60$ von $100$ Kästchen $= \frac{60}{100}$}{}
\end{pfloesung}
\begin{pfloesung}{}
\lz{2b)}{$40\,\%$}{$\frac{2}{5} = \frac{40}{100}$}{}
\lz{c)}{$35\,\%$}{$\frac{7}{20} = \frac{35}{100}$}{}
\end{pfloesung}
\begin{pfloesung}{}
\lz{3b)}{$28\,\%$}{mal $4$ erweitern: $\frac{28}{100}$}{}
\lz{c)}{$18\,\%$}{mal $2$ erweitern: $\frac{18}{100}$}{}
\lz{d)}{$65\,\%$}{mal $5$ erweitern: $\frac{65}{100}$}{}
\lz{e)}{$75\,\%$}{mal $25$ erweitern: $\frac{75}{100}$}{}
\end{pfloesung}
\begin{pfloesung}{}
\lz{4.}{b) $\frac{60}{100}$; $60\,\%$; $0{,}6$}{$\frac{3}{5} = \frac{60}{100}$}{}
\lz{}{c) $\frac{13}{20}$; $\frac{65}{100}$; $65\,\%$}{$0{,}65 = \frac{65}{100}$, gekürzt mit $5$}{}
\lz{}{d) $\frac{1}{25}$; $\frac{4}{100}$; $0{,}04$}{$4\,\% = \frac{4}{100}$, gekürzt mit $4$}{}
\lz{}{e) $\approx 33{,}3\,\%$; $\approx 0{,}333$}{$1 : 3 = 0{,}333\ldots$; mit Nenner $100$ geht nicht genau}{}
\end{pfloesung}
\begin{pfloesung}{}
\lz{5a)}{jedes fünfte Kästchen ($20\,\%$)}{$20$ von $100$ grau}{}
\lz{b)}{$5$ von $100$ ($5\,\%$)}{$5$ Kästchen grau}{}
\end{pfloesung}
\begin{pfloesung}{}
\lz{6b)}{$10\,\%$}{$\frac{1}{10} = \frac{10}{100}$}{}
\lz{c)}{$2\,\%$}{$\frac{1}{50} = \frac{2}{100}$}{}
\lz{d)}{$60\,\%$}{$\frac{3}{5} = \frac{60}{100}$}{}
\lz{e)}{jeder zweite ($2.$)}{$50\,\% = \frac{1}{2}$, jeder zweite}{}
\lz{f)}{jeder zwanzigste ($20.$)}{$5\,\% = \frac{5}{100} = \frac{1}{20}$, jeder zwanzigste}{}
\lz{g)}{$8$ von $100$}{$8\,\% = 8$ von $100$}{}
\end{pfloesung}
\begin{pfloesung}{}
\lz{7.}{$5$ von $100$ Paketen kommen zu spät.}{$5\,\% = \frac{5}{100}$; jedes 5. wären $20\,\%$}{}
\end{pfloesung}
\begin{pfloesung}{}
\lz{8a)}{$19\,\%$}{$100 - 35 - 28 - 10 - 8 = 19$}{}
\lz{b)}{Milch}{jedes zehnte $= \frac{1}{10} = 10\,\%$}{}
\end{pfloesung}
\begin{pfloesung}{}
\lz{9.}{„Jeder zehnte Grundschüler …“ ist falsch. Richtig: Jeder zwanzigste Grundschüler frühstückt nie zu Hause (auch richtig: $5$ von $100$ Grundschülern …; jeder zehnte Oberschüler …).}{$5\,\% = \frac{5}{100} = \frac{1}{20}$; jeder zehnte ($10\,\%$) gilt für die Oberschule; $\frac{3}{5} = 60\,\%$ stimmt}{}
\end{pfloesung}
\end{document}
